\chapter{Resultados}
\label{ch:resultados}
% Énfasis Ciencia de Datos: toda afirmación con tabla/figura. Describe; la interpretación vive en Conclusiones.

Este capítulo reúne lo que la matriz de experimentos produjo y lo lee con el protocolo de la Sección~\ref{sec:protocolo-analisis}. Las cuatro primeras secciones comprueban que los datos sirven para la pregunta, y cada una mira un solo lado de la correlación. La Sección~\ref{sec:res-cobertura} cuenta qué variables quedan medidas y cuánta censura deja la velocidad. La Sección~\ref{sec:res-rango} mide qué parte de cada predictor mueve el \emph{learning rate} y qué parte la \emph{seed}. La Sección~\ref{sec:res-solape} acota en qué ventanas la velocidad sigue por delante de la medición. La Sección~\ref{sec:res-val-test} comprueba que la curva de validación cuenta lo mismo que el test. Las secciones siguientes responden a las seis hipótesis en el orden de los objetivos. La Sección~\ref{sec:res-por-condicion} responde a H1 y la Sección~\ref{sec:res-incremental} a H2, seguida del coste de cada predictor frente a su capacidad predictiva en la Sección~\ref{sec:res-coste}. La Sección~\ref{sec:res-familia} responde a H3, la Sección~\ref{sec:res-temporal} a H4, la Sección~\ref{sec:res-robustez} a H5 y la Sección~\ref{sec:res-concordancia} a H6. Cada una describe lo que sale y cierra con el veredicto de su hipótesis. La interpretación y el contraste con la literatura quedan para el Capítulo~\ref{ch:conclusiones}. La Sección~\ref{sec:res-conclusiones} recoge los veredictos.

\section{Qué variables quedan medidas}
\label{sec:res-cobertura}

Seis de las siete variables dependientes se leen sin compararlas con ningún umbral, y la matriz las entrega casi completas, entre 919 y 921 valores utilizables de 960. Lo que falta viene siempre del mismo sitio. En 41 entrenamientos el \emph{loss} pasó a NaN y no se recuperó, todos con SGD, 28 con la red \emph{fully connected} y 13 con la convolucional. Ninguna ResNet-18 divergió, y la Sección~\ref{sec:setup} atribuye esa diferencia a su \emph{batch normalization}. La celda peor servida conserva 30 de sus 40 entrenamientos.

El \emph{accuracy} de test y el \emph{gap} de \emph{accuracy} tienen valor en los 960 entrenamientos, pero en esos 41 el valor no es una medición. Una red divergida termina con todos los pesos en NaN. Sus salidas son NaN, y el \emph{argmax} sobre ellas devuelve siempre la clase cero. La cifra registrada es la frecuencia de esa clase en el conjunto de test. Las variables que se leen del \emph{loss} no tienen ese problema, porque el NaN se propaga por la aritmética y el valor queda ausente. Esos 41 valores se cuentan aparte y no entran en ninguna variable leída al final del entrenamiento. Dos de los 41, la red \emph{fully connected} sobre MNIST con SGD y el \emph{learning rate} 0,3 en dos \emph{seeds}, aprendieron en la primera \emph{epoch} y divergieron en la segunda. Por eso cuentan como entrenamientos que aprendieron y entran censurados en la velocidad, según fija la Sección~\ref{sec:protocolo-poblacion}.

La séptima variable, las \emph{epochs} hasta el umbral, es la única que se censura. La tienen definida 573 de los 960 entrenamientos, entre 15 y 38 por celda, y 22 de las 24 celdas llegan a 20. Contada en entrenamientos, la pérdida parece grave. Contada en la unidad que usa el estadístico, lo es menos. Una correlación de rangos no consume entrenamientos, consume pares comparables, parejas de las que se sabe cuál fue más rápida. Un entrenamiento censurado no cruzó dentro del presupuesto. Fue más lento que cualquiera que sí cruzó, y forma par con todos ellos. Solo se pierden los pares de dos censurados, porque entre ellos no se sabe cuál habría cruzado antes. Medida en pares, cada celda conserva entre el 62 y el 99\,\% de los suyos, con mediana del 82\,\% y 10 celdas por encima del 85\,\%. La más pobre es Tiny-ImageNet con la red \emph{fully connected} y Adam, que cruza 15 veces y aun así conserva el 62\,\% de sus pares. Esa fracción es una función del número de cruces, y es la misma pérdida expresada en la escala en la que trabaja el estadístico.

Los entrenamientos censurados tampoco se reparten al azar por la rejilla. La Figura~\ref{fig:ventana-lr} recorre los ocho \emph{learning rates} de cada optimizador y muestra cuántos de los cinco entrenamientos de cada uno cruzan el umbral. Los dos paneles comparten el eje, y una misma columna es el mismo valor de \emph{learning rate} en los dos, aunque cada lista empiece en un sitio distinto. Cada fila lleva su umbral, tomado de la Tabla~\ref{tab:umbrales}, porque de él dependen la posición y la anchura de su tramo. En las 24 celdas los cruces ocupan un tramo contiguo, sin un solo hueco, y ese tramo se desplaza hacia los valores mayores a medida que la red gana profundidad. La red \emph{fully connected} cruza en las cinco primeras posiciones de su lista y casi nunca después, la convolucional casi siempre entre la segunda y la sexta, y la ResNet-18 es la única que llega a las dos últimas. Los dos optimizadores muestran el mismo desplazamiento. Ese desplazamiento no se debe solo a la arquitectura, porque el umbral también sube con la profundidad y a la red más honda se le pide más \emph{accuracy}. Las dos causas empujan en el mismo sentido, y esta figura no las separa.

\begin{figure}[ht]
\centering
\includegraphics{ventana-lr}
\caption{Entrenamientos que cruzan el umbral en cada \emph{learning rate}, por celda.}
\label{fig:ventana-lr}
\end{figure}

La censura resulta así una consecuencia del diseño. Los extremos de la rejilla existen para producir entrenamientos malos, como expone la Sección~\ref{sec:matriz}, y son precisamente ellos los que no alcanzan el umbral. Descartar los entrenamientos censurados dejaría solo el centro de cada rejilla, y retiraría de la muestra la variación que la matriz se construyó para crear.

El nivel alcanzado cuadra con lo esperado. Tomando en cada celda el \emph{learning rate} cuya mediana de \emph{accuracy} de test sobre las cinco \emph{seeds} es mayor, nueve de las 12 combinaciones de conjunto y arquitectura caen dentro de la banda de la Tabla~\ref{tab:referencias}. La red \emph{fully connected} sobre CIFAR-10 llega al 57,8\,\%, por encima del 50 al 54\,\% esperado. La convolucional sobre CIFAR-100 se queda en el 39,0\,\%, por debajo del 45 al 55\,\%, y la ResNet-18 sobre CIFAR-10 en el 83,6\,\%, justo por debajo del 84 al 92\,\%.

\section{Rango dinámico del predictor}
\label{sec:res-rango}

Un coeficiente de correlación nulo admite tres lecturas que desde fuera no se distinguen. Puede que la métrica no prediga nada. Puede que el suceso a predecir ya hubiera ocurrido cuando se midió la métrica. O puede que la métrica no se mueva dentro de una celda más de lo que la mueve el ruido de la \emph{seed}. Esta sección descarta la tercera. Mide qué parte de la dispersión de cada columna registrada explica el \emph{learning rate} y qué parte explica la \emph{seed}. Es una comprobación del lado del predictor, sin consultar ninguna variable dependiente.

Dentro de una celda hay 40 entrenamientos, ocho \emph{learning rates} por cinco \emph{seeds}, y dentro de cada \emph{learning rate} lo único que cambia es la \emph{seed}. Los valores de una columna se sustituyen por su posición en la ordenación, del 1 al 40. La dispersión de esas posiciones se parte en dos, la que separa a los grupos entre sí y la que queda dentro de cada grupo:

\begin{equation}
\eta^2_{\mathrm{LR}} = \frac{\sum_{g} n_g\,(\bar{R}_g - \bar{R})^2}{\sum_{i} (R_i - \bar{R})^2},
\label{eq:eta-rangos}
\end{equation}

donde $R_i$ es la posición del entrenamiento $i$, $\bar{R}$ la media de todas las posiciones, $\bar{R}_g$ la media del grupo $g$ y $n_g$ su tamaño. El cociente vale 1 cuando el \emph{learning rate} explica toda la dispersión y 0 cuando las medias de los grupos coinciden. El mismo cociente se calcula agrupando por \emph{seed} en lugar de por \emph{learning rate}, y de ahí salen las dos cifras que se comparan. Trabajar sobre posiciones impide que un entrenamiento extremo domine el reparto, y hace la cifra comparable entre métricas cuyas unidades difieren en órdenes de magnitud.

El cociente no vale cero aunque el \emph{learning rate} no intervenga. Si se reparten $n$ posiciones al azar entre $k$ grupos, las medias de los grupos difieren algo por pura suerte, y el valor esperado del cociente es $(k-1)/(n-1)$. Esa es la referencia contra la que se lee cada cifra. Vale para cualquier conjunto fijo de valores, y trabajar con posiciones hace además que no dependa de cómo se distribuye la columna.

El estadístico no exige normalidad ni varianzas iguales, porque opera sobre posiciones y no contrasta ninguna hipótesis. Tampoco exige que la relación con el \emph{learning rate} sea monótona, porque trata los ocho valores como grupos sin orden. Lo que sí exige es la rejilla cruzada y completa de la Sección~\ref{sec:matriz}, que garantiza que dentro de un grupo solo cambie la \emph{seed}. Cada pareja de \emph{learning rate} y \emph{seed} tiene un solo entrenamiento, y no se puede saber si una \emph{seed} sale alta en todos los \emph{learning rates} o solo en algunos. Por eso las dos cifras se publican aparte y no como un reparto que sume uno.

La Figura~\ref{fig:rango-celda} recorre una celda completa, MNIST con la red convolucional y SGD, con la \emph{gradient disparity} medida al 5\,\% del presupuesto. Las medias de los ocho grupos se reparten por toda la escala. Dentro de un grupo, los cinco entrenamientos quedan en una banda más estrecha, que se ensancha donde la columna vale más. El panel derecho cambia cada valor por su posición dentro de la celda, de 1 el más bajo a 39 el más alto, que es la vista sobre la que opera la Ecuación~\ref{eq:eta-rangos}. Son 39 porque el entrenamiento divergido de la celda no deja valor en esta columna. Entre los 39 están los 12 que se quedaron en el azar, y la figura los marca aparte porque son parte de lo que hay que ver. El \emph{learning rate} explica 0,88 de la dispersión, frente a una referencia de 0,18 para sus ocho grupos. La \emph{seed} explica 0,08, por debajo de la referencia de 0,11 de sus cinco grupos, es decir, nada más allá del azar.

\begin{figure}[ht]
\centering
\includegraphics{rango-celda}
\caption{Una celda completa, MNIST con la red convolucional y SGD, al 5\,\%.}
\label{fig:rango-celda}
\end{figure}

El panel izquierdo muestra además una forma que el cociente no recoge. La columna sube hasta el \emph{learning rate} $10^{-3}$ y a partir de ahí baja, sin crecer ni decrecer de un extremo al otro de la rejilla. El cociente la puntúa alto igual, porque mide si hay estructura, no en qué sentido va. Un coeficiente de rangos, en cambio, solo detecta relaciones que suben o bajan en un único sentido, y leería la mitad de esa curva como ausencia de señal. Dentro de una celda, la relación entre una métrica y una variable es la composición de dos curvas sobre el \emph{learning rate}, y ningún estadístico de rangos la deshace. La Sección~\ref{sec:protocolo-analisis} declara las celdas donde esa relación no puede ser monótona y añade una lectura, la de selección, que no depende de la forma.

Cada métrica entra representada por su columna titular, la que fija la Tabla~\ref{tab:columnas-titulares}, porque las demás columnas de una métrica son otras lecturas de la misma cantidad (Sección~\ref{sec:variables}). La Figura~\ref{fig:rango-columnas} repite la cuenta en las 24 celdas al 5\,\% del presupuesto y resume cada columna por su mediana. Las ocho métricas de gradiente superan su referencia con holgura, desde 0,45 la más floja hasta 0,92 la que más. La cuenta se extiende a las cuatro ventanas tempranas. De las 768 casillas de métrica de gradiente, que son 24 celdas por cuatro ventanas por ocho columnas, solo 30 quedan en la referencia o por debajo, un 3,9\,\%. Todas corresponden a la familia de variabilidad o a la m-coherencia, que da exactamente el mismo número que la escala de ruido por la Ecuación~\ref{eq:gns-mcoh}. La tercera lectura de un coeficiente nulo queda descartada en casi toda la rejilla.

\begin{figure}[ht]
\centering
\includegraphics{rango-columnas}
\caption{Dispersión explicada por el \emph{learning rate} en las 24 celdas.}
\label{fig:rango-columnas}
\end{figure}

La parte que queda dentro de los \emph{learning rates} es el complemento de esa cifra, y es la única que puede usar la cuenta del contraste que fija el \emph{learning rate}, la granulada de la Sección~\ref{sec:protocolo-granulada}. Al 5\,\% es el 8\,\% de la dispersión en la \emph{stiffness}, el 10\,\% en la \emph{gradient disparity}, el 22\,\% en la \emph{gradient confusion} y el 33\,\% en GWA. En GSNR es el 46\,\% y en la escala de ruido el 55\,\%. En las tres lecturas del predictor de referencia (\emph{baseline}) queda entre el 3 y el 5\,\%. Las cuentas sobre las 24 celdas usan solo los entrenamientos que aprendieron, según el criterio de la Sección~\ref{sec:calibracion}, mientras que la celda de ejemplo muestra todos los que tienen valor. Repetir las 24 celdas al 5\,\% con todos los que tienen valor mueve las medianas menos de tres centésimas. Esa coincidencia no aísla el efecto de los entrenamientos que nunca aprendieron, porque quitarlos recorta también el extremo alto de la rejilla. Dice solo que el resultado no depende de qué población se elija.

Las tres columnas que encabezan la Figura~\ref{fig:rango-columnas} no cuestan nada de calcular. TSE marca 0,97 y las dos lecturas de la curva de validación 0,95, por encima del 0,92 de la métrica instrumentada que más se mueve. Dentro de una celda, la columna que mejor sigue al \emph{learning rate} es justamente la que no hay que pagar, y esa es la vara que la hipótesis H2 obliga a superar.

Ese mismo resultado deja medida la causa común que la Sección~\ref{sec:matriz} declara. Dentro de una celda, la métrica y el indicador de eficiencia responden los dos al \emph{learning rate}, y una correlación entre ellos mide en buena parte si la métrica es una relectura del \emph{learning rate} en el mismo sentido que el resultado. Por eso la posición en la rejilla entra en todas las tablas del contraste como predictor de coste cero, y la cuenta granulada de la Sección~\ref{sec:protocolo-granulada}, que solo compara entrenamientos con el mismo \emph{learning rate}, acompaña a cada recuento.

\section{Solape entre la ventana y el desenlace}
\label{sec:res-solape}

De las tres lecturas de un coeficiente nulo con las que abre la Sección~\ref{sec:res-rango} queda por acotar la segunda, que el suceso a predecir ya hubiera ocurrido al medir. La Sección~\ref{sec:ventana} lo enuncia como condición y la Sección~\ref{sec:riesgos} exige medir su alcance antes de interpretar cualquier resultado de velocidad. Esta sección lo mide celda a celda y ventana a ventana sobre las tres variables que se leen a lo largo de la curva de validación, las dos de velocidad y el mejor \emph{loss}, y dice qué ventanas conservan sentido para cada una. Como la anterior, trabaja sobre un solo lado, aquí el de las variables dependientes, y no consulta ninguna métrica.

Una ventana se cierra en la \emph{epoch} en la que se lee el predictor, la 1, 2, 5 y 10 en MNIST y la 2, 4, 10 y 20 en los otros tres conjuntos. Un suceso que ocurre en esa misma \emph{epoch} cuenta como ya pasado, porque la métrica y el suceso describen la misma red al final de la misma \emph{epoch}, y la primera no puede anticipar lo que ya es. El suceso de las \emph{epochs} hasta el umbral es el cruce. El del mejor \emph{loss} de validación es la \emph{epoch} en la que la curva suavizada toca su mínimo. El área bajo la curva no tiene suceso, porque es una suma ponderada de los \emph{loss} de todas las \emph{epochs}. Lo que la ventana deja atrás es la parte de esa suma que ya está fijada al cerrarse.

La unidad de la cuenta son los pares comparables, la misma con la que la Sección~\ref{sec:res-cobertura} mide la censura, porque un coeficiente de rangos consume pares. Un par sigue siendo una predicción cuando ninguno de sus dos entrenamientos ha tenido el suceso al cerrar la ventana, y un entrenamiento censurado nunca lo tiene. Con $k$ cruces en la celda, $f$ de ellos todavía por delante y $c$ entrenamientos censurados, la parte de los pares comparables que sigue por delante es
\begin{equation}
\frac{\binom{f}{2} + f\,c}{\binom{k}{2} + k\,c},
\label{eq:pares-delante}
\end{equation}
y para el mejor \emph{loss}, que nunca se censura, $c$ vale cero. Para el área, lo que queda por delante de cada entrenamiento es el complemento de la fracción ya fijada, y la celda se resume por su mediana.

Una ventana sirve para una variable en una celda cuando al menos la mitad de esa variable sigue por delante, porque entonces la ventana predice para la mayor parte de la información que usa el contraste. La regla se fijó después de conocer las cifras agregadas sobre los 960 entrenamientos y antes del desglose por celda, y se declara como lo que es, una regla posterior a los datos (Sección~\ref{sec:protocolo-cuando}). Desde la revisión del método la regla ya no decide qué celdas se leen, solo describe cuánto queda por delante en cada una. La velocidad se lee por hitos, entre los entrenamientos que aún no han cruzado al cerrar la ventana, según fija la Sección~\ref{sec:protocolo-hitos}.

La Figura~\ref{fig:solape-celda} recorre CIFAR-10 con la red convolucional y SGD, la celda que reparte sus cruces cinco a cinco entre las cuatro ventanas. Cada curva es el \emph{accuracy} de validación suavizado de un entrenamiento, la línea horizontal es el umbral y las cuatro verticales marcan la \emph{epoch} en la que cierra cada ventana. De sus 40 entrenamientos cruzan 21, y las cuatro ventanas dejan atrás cinco, 10, 15 y 20 de esos cruces. El panel inferior los cuenta acumulados, porque en las curvas se amontonan a la izquierda y el reparto no se ve. En pares, la parte que sigue por delante baja de 0,70 al 5\,\% a 0,43 al 10\,\%, 0,21 al 25\,\% y 0,03 al 50\,\%. A partir de la segunda ventana, lo que una métrica leída ahí puede hacer con la variable principal es, en su mayor parte, describir un cruce que ya ha pasado.

\begin{figure}[ht]
\centering
\includegraphics{solape-celda}
\caption{Cruces del umbral en CIFAR-10 con la red convolucional y SGD.}
\label{fig:solape-celda}
\end{figure}

Contado en entrenamientos y sumado sobre las 24 celdas, de los 573 cruces ya ha ocurrido el 10\,\% al cerrar la ventana del 5\,\% y el 35\,\% al cerrar la del 10\,\%. Al cerrar la del 25\,\% ha ocurrido el 72\,\%, y al cerrar la del 50\,\%, el 85\,\%. Esta cuenta va en entrenamientos, no en pares, y mide el tamaño del solape sin decidir nada.

El reparto entre conjuntos de datos confunde si se lee en fracción del presupuesto. La \emph{epoch} mediana de cruce es la 5 en MNIST, en CIFAR-10 y en Tiny-ImageNet, y la 4 en CIFAR-100. Los cuatro conjuntos cruzan a la misma altura absoluta. Lo que cambia es el presupuesto, 20 \emph{epochs} en MNIST frente a 40 en los demás, y la misma \emph{epoch} es en MNIST el doble de fracción. Las seis celdas de MNIST conservan la mitad por delante al 10\,\% porque su ventana del 10\,\% cierra en la \emph{epoch} 2, y la de los otros conjuntos en la 4. Seis de las 12 celdas que sirven al 10\,\% lo son por ese denominador. Por arquitectura, la red \emph{fully connected} cruza antes que las otras dos, con mediana en la \emph{epoch} 4 frente a 5 en la ResNet-18 y 6 en la convolucional. Tampoco ahí se separa la dificultad de la exigencia, porque a la red plana se le pide el umbral más bajo de las tres. Comparar el solape entre conjuntos mezcla tiempo absoluto y relativo. Dentro de una celda es constante.

Para las \emph{epochs} hasta el umbral, contadas celda a celda y en pares, la ventana del 5\,\% sirve en 22 celdas y la del 10\,\% en 12, y ninguna de las dos posteriores sirve en ninguna. Las dos celdas que fallan ya al 5\,\% son la red \emph{fully connected} con Adam en CIFAR-100 y en Tiny-ImageNet, donde más de la mitad de los cruces, 11 de 20 y 10 de 15, caen en las dos primeras \emph{epochs}. Al 10\,\% el reparto se parte por arquitectura. Quedan por encima de la mitad las seis celdas de MNIST, las cuatro de la convolucional y la ResNet-18 sobre Tiny-ImageNet y las dos de la ResNet-18 con SGD sobre CIFAR. Ninguna red \emph{fully connected} fuera de MNIST llega a la mitad, y la que más conserva, Tiny-ImageNet con SGD, se queda en 0,42. Es el mismo orden por profundidad que dibuja la Figura~\ref{fig:ventana-lr}, porque a la red plana se le pide el umbral más bajo y lo alcanza antes.

El mejor \emph{loss} de validación conserva más ventanas, 24 celdas al 5\,\%, 17 al 10\,\%, nueve al 25\,\% y dos al 50\,\%, y lo que lo consume es el sobreajuste. Sin \emph{data augmentation} ni \emph{weight decay}, la curva de \emph{loss} de validación toca fondo pronto y vuelve a subir (Sección~\ref{sec:setup}). En CIFAR-10 con ResNet-18 y Adam, 37 de los 40 entrenamientos tocan fondo en la \emph{epoch} 10 o antes y el último en la 14, y al 25\,\% queda por delante el 0,4\,\% de los pares. En esas celdas el suceso de esta variable es el inicio del sobreajuste, y así hay que leerla.

El comportamiento del área bajo la curva es casi mecánico. Al cerrar cada ventana queda fijada casi exactamente la fracción del presupuesto transcurrida, con medianas de 0,96, 0,90, 0,76 y 0,51 por delante, frente a los 0,96, 0,91, 0,76 y 0,50 que dejaría una curva plana de 40 \emph{epochs}. Lo dictan los pesos de la suma, sin que los datos aporten nada. La variable sirve hasta el 25\,\% en las 24 celdas. Al 50\,\% cae sobre la línea por construcción, con 13 celdas de 24, y para esa ventana queda declarada como descrita a medias.

Las cuentas anteriores usan los 960 entrenamientos. Repetidas solo con los que aprendieron, según el criterio de la Sección~\ref{sec:calibracion}, el recuento de celdas del área no cambia. El de las \emph{epochs} hasta el umbral pierde una celda al 5\,\%, y el del mejor \emph{loss} se mueve una celda al 5, al 10 y al 50\,\%. Un entrenamiento que nunca aprende nunca cruza, entra como censurado, y quitarlo apenas toca el reparto.

Para las variables de velocidad, el barrido de ventanas de la hipótesis H4 queda acotado a las del 5 y el 10\,\%. Se declara en \emph{epochs}, la 1 frente a la 2 en MNIST y la 2 frente a la 4 en los demás conjuntos. Las tres variables que se miden al terminar el entrenamiento no tienen este límite. La cuenta es además un suelo. Un entrenamiento al 95\,\% del umbral al cerrar la ventana tiene el desenlace casi escrito y aquí cuenta como por delante. El solape real es mayor que el medido. La Sección~\ref{sec:protocolo-hitos} fija qué hace el contraste con los entrenamientos que ya han cruzado, que es dejarlos fuera de la lectura primaria de la velocidad. La cuenta de esta sección solo describe cuánto queda por delante en cada celda. El solape favorece además al predictor de referencia. En un entrenamiento que ya cruzó, el \emph{accuracy} de validación leído en la ventana vale ya el umbral, y lo detecta sin coste. Por eso H2 no se decide sobre la velocidad, según fija la Sección~\ref{sec:protocolo-eval}.

\section{Concordancia entre validación y test}
\label{sec:res-val-test}

Las variables de velocidad se leen sobre la curva de validación, y el test se evalúa una sola vez al terminar, como fija la Sección~\ref{sec:protocolo-eval}. Lo que este trabajo afirme sobre velocidad descansa por tanto en que la validación cuente lo mismo que el test. El diseño lo hace esperable, porque las dos particiones salen del mismo origen y ninguna interviene en el entrenamiento. Esta sección lo comprueba en el único instante en que las dos lecturas coexisten, el final del entrenamiento, sobre los 919 entrenamientos que no divergieron. Es una correlación entre dos variables dependientes y tampoco consulta ninguna métrica.

Se comparan dos cosas, el orden y el nivel. El orden importa porque los contrastes de la Sección~\ref{sec:protocolo-analisis} ordenan los entrenamientos de una celda. El nivel importa porque el umbral de la Sección~\ref{sec:calibracion} se fijó sobre validación. El orden se mide con la tau de Kendall dentro de cada celda, que cuenta los pares de entrenamientos que las dos lecturas ordenan igual, resta los que ordenan al revés y divide por el total de pares. Su mediana sobre las seis celdas de cada conjunto es 0,93 en los dos CIFAR, 0,92 en Tiny-ImageNet y 0,84 en MNIST para el \emph{accuracy}, y entre 0,86 y 0,94 para el \emph{loss}. Las 24 celdas caen entre 0,71 y 0,97. El mínimo es MNIST con ResNet-18 y Adam, cuyos 40 entrenamientos terminan en un tramo de 0,008 de \emph{accuracy} de test, del tamaño del ruido de medida, y con esa separación el orden es ruido en las dos lecturas a la vez. Lo mismo ocurre en cinco de las seis celdas de MNIST entre los entrenamientos que aprendieron y no divergieron, que terminan a menos de cinco centésimas unos de otros. Es un aviso para los contrastes sobre el rendimiento final, porque en MNIST esa variable apenas tiene recorrido entre los entrenamientos sanos. El nivel es la diferencia validación menos test de cada entrenamiento, leída contra el ruido de medir el \emph{accuracy} sobre un número finito de ejemplos. Su mediana es 0,003 en CIFAR-10 y cero en los otros tres conjuntos, y el 95\,\% de los entrenamientos queda a menos de 0,011 de la igualdad. Los que se separan más de dos desviaciones típicas del ruido van del 0\,\% en Tiny-ImageNet al 1,8\,\% en CIFAR-10, por debajo del 4,6\,\% que el ruido solo produciría, porque los entrenamientos clavados en el azar dan el mismo \emph{accuracy} en las dos particiones.

La misma evaluación cierra la comprobación anunciada en la Sección~\ref{sec:protocolo-eval}. La diferencia entre la F1 macro y el \emph{accuracy} de test tiene mediana 0,002. Entre los 804 entrenamientos que aprendieron y no divergieron, el 95\,\% queda por debajo de 0,018 y el máximo es 0,071. Repetidas solo con los entrenamientos que aprendieron, las medianas de la tau bajan menos de tres centésimas, porque desaparecen los empates del azar, y el nivel no cambia. La curva de validación es una lectura fiel del test en orden y en nivel, y las variables de velocidad que se leen sobre ella no son un artefacto de la partición. Lo que no se puede establecer es el acuerdo en cada \emph{epoch}, porque el test se mira una sola vez.

\section{Existencia de la asociación}
\label{sec:res-por-condicion}
% Escrita el 2026-09-03 (fase B, paso 4) y reescrita esa misma tarde con el método revisado.
% Números de results/tabla_larga.parquet, regenerable con `uv run python src/contrast.py`;
% figura signos-h1 en src/figures.py. H1 se lee como requisito previo, no como hallazgo. La lectura de
% selección (results/seleccion.parquet) es de H2 y va en §Valor incremental.

Esta sección responde a la hipótesis H1, la de existencia, sobre la familia primaria que fija la Sección~\ref{sec:protocolo-agregacion}. La familia es la ventana del 5\,\%, tres variables y diez predictores. Las variables son las \emph{epochs} hasta el umbral leídas por hitos (\emph{landmark}), es decir, solo entre los entrenamientos que aún no habían cruzado al cerrar la ventana, el \emph{accuracy} de test y el \emph{gap} de \emph{loss}. Los predictores son la posición en la rejilla, las tres lecturas del predictor de referencia y las seis métricas de gradiente que quedan tras la poda. Cada D se calcula dentro de una celda, sobre los entrenamientos que aprendieron. D vale 1 cuando el predictor ordena los entrenamientos igual que la variable, $-1$ cuando los ordena al revés y 0 cuando acierta tantos pares como falla. Cada D lleva su intervalo por \emph{jackknife}, y una celda muestra asociación cuando ese intervalo deja fuera el cero. La Tabla~\ref{tab:signos-h1} da, para cada combinación, cuántas celdas tienen D positiva y cuántas negativa, cuántas de cada signo muestran asociación, y el mismo recuento de signos para la cuenta granulada, la que compara solo entrenamientos con el mismo \emph{learning rate}. En el texto, el número de celdas de un predictor es el de las que muestran asociación con su signo mayoritario, el signo que su D tiene en más celdas; la tabla da también las del signo contrario. La Figura~\ref{fig:signos-h1} resume las 24 D de cada combinación en una caja. La raya es la mediana, la caja abarca la mitad central de las celdas y los bigotes van de la celda más baja a la más alta. Las cajas de las tres lecturas del predictor de referencia quedan lejos del cero con la velocidad y con el \emph{accuracy} de test, y con el \emph{gap} cruzan el cero y son anchas. Las cajas de las seis métricas cruzan el cero, con tres excepciones, GWA con la velocidad y con el \emph{accuracy} de test y la \emph{stiffness} con el \emph{gap}.

\begin{figure}[ht]
\centering
\includegraphics{signos}
\caption{D por celda de la familia primaria al 5\,\%.}
\label{fig:signos-h1}
\end{figure}

\begin{table}[ht]
\centering
\footnotesize
\setlength{\tabcolsep}{3.5pt}
\caption{Recuento de celdas de la familia primaria. En negrita, las combinaciones con asociación del mismo signo en 15 celdas o más.}
\label{tab:signos-h1}
\begin{tabular}{lccccccccc}
\toprule
& \multicolumn{3}{c}{\textbf{\emph{Epochs} hasta el umbral}} & \multicolumn{3}{c}{\textbf{\emph{Accuracy} de test}} & \multicolumn{3}{c}{\textbf{\emph{Gap} de \emph{loss}}} \\
\cmidrule(lr){2-4}\cmidrule(lr){5-7}\cmidrule(lr){8-10}
\textbf{Predictor} & $+/-$ & asoc. & gran. & $+/-$ & asoc. & gran. & $+/-$ & asoc. & gran. \\
\midrule
Posición en la rejilla & 9/15 & 5/7 & 0/0 & 11/13 & 6/7 & 0/0 & 21/3 & \textbf{15/0} & 0/0 \\
\emph{Loss} de validación & 24/0 & \textbf{24/0} & 16/6 & 0/24 & \textbf{0/19} & 4/18 & 7/17 & \textbf{6/16} & 7/16 \\
\emph{Accuracy} de validación & 0/24 & \textbf{0/24} & 5/18 & 23/1 & \textbf{18/0} & 19/5 & 17/7 & \textbf{16/6} & 15/8 \\
TSE & 24/0 & \textbf{23/0} & 16/6 & 1/23 & \textbf{0/17} & 7/17 & 8/16 & \textbf{7/16} & 11/12 \\
\addlinespace
GNS & 14/10 & 7/3 & 11/12 & 14/10 & 3/6 & 9/12 & 9/15 & 5/6 & 15/9 \\
GSNR & 11/13 & 6/7 & 11/11 & 15/9 & 7/4 & 13/9 & 8/16 & 2/11 & 6/16 \\
\emph{Gradient disparity} & 9/15 & 4/9 & 9/12 & 15/9 & 12/4 & 13/10 & 9/15 & 4/11 & 14/7 \\
\addlinespace
\emph{Stiffness} & 15/8 & 9/5 & 14/8 & 12/12 & 5/6 & 13/11 & 4/20 & \textbf{1/15} & 10/12 \\
\emph{Gradient confusion} & 12/12 & 5/4 & 10/10 & 11/13 & 5/5 & 18/6 & 14/10 & 9/5 & 7/13 \\
GWA & 21/3 & \textbf{16/1} & 9/13 & 4/20 & \textbf{1/15} & 7/16 & 8/16 & 6/4 & 11/11 \\
\bottomrule
\end{tabular}
\end{table}

Antes de leer los signos, el protocolo pide declarar tres cifras de la población. En los contrastes del \emph{gap} entran solo los entrenamientos que alcanzaron el umbral de su celda sobre el conjunto de entrenamiento, y son 693 de los 806, con 25 por celda de mediana. Cinco celdas conservan los 40 y tres se quedan en 20, Tiny-ImageNet con la red \emph{fully connected} y con la convolucional bajo Adam, y MNIST con la red \emph{fully connected} bajo SGD. Al cerrar la ventana del 5\,\% quedan por cruzar entre 15 y 40 entrenamientos por celda, la mitad de las celdas con 30 o más, y entre 50 y 713 pares comparables, 366 de mediana. Las dos celdas con menos entrenamientos en riesgo son la red \emph{fully connected} bajo Adam sobre Tiny-ImageNet, con 15, y sobre CIFAR-100, con 17, las mismas que la Sección~\ref{sec:res-solape} señala. Queda la forma de las curvas. El protocolo advierte de que D no puede ver una relación cuando el predictor solo sube o solo baja a lo largo del \emph{learning rate} y la variable pasa por un pico o un valle. Esas celdas son pocas para todos los predictores. La \emph{stiffness} es la peor, con seis en velocidad, siete en \emph{accuracy} de test y cinco en el \emph{gap}. La escala de ruido, GSNR y la \emph{gradient confusion} tienen entre dos y cuatro, y la disparidad y GWA entre dos y tres.

La posición en la rejilla es el logaritmo del \emph{learning rate}, y marca el techo de lo que una lectura monótona del \emph{learning rate} puede dar. Con la velocidad muestra asociación en 12 celdas en total y con el \emph{accuracy} de test en 13, pero repartidas casi por igual entre los dos signos, cinco positivas y siete negativas en la primera, seis y siete en la segunda. Con su signo mayoritario son siete y siete. Las dos variables pasan por un valle o por un pico a lo largo de la rejilla, y el signo depende de a qué lado del óptimo caigan los \emph{learning rates} de la celda. Con el \emph{gap} de \emph{loss} la relación sí es monótona. Un \emph{learning rate} mayor acompaña a un \emph{gap} mayor, y la posición muestra asociación positiva en 15 celdas y negativa en ninguna, cinco por arquitectura.

Las tres lecturas del predictor de referencia ordenan los entrenamientos de casi todas las celdas. El \emph{accuracy} de validación muestra asociación negativa con las \emph{epochs} hasta el umbral en las 24 celdas, y el \emph{loss} de validación y TSE positiva en 24 y en 23, con medianas de 0,64 a 0,71 en valor absoluto. Con el \emph{accuracy} de test el signo se invierte, como corresponde, en 18, 19 y 17 celdas, con medianas de 0,33 a 0,40 y sin ninguna celda del signo contrario. Con el \emph{gap} de \emph{loss} los tres se parten, y el corte es la arquitectura. En las ocho celdas de la red \emph{fully connected} y en las ocho de la convolucional, un \emph{loss} de validación alto en la ventana acompaña a un \emph{gap} menor al terminar, con asociación en las 16. En la ResNet-18 acompaña a un \emph{gap} mayor, con asociación en seis de las ocho. El \emph{accuracy} de validación dibuja lo mismo con el signo opuesto. Contadas las 24 celdas juntas, las dos relaciones se anulan.

De las seis métricas de gradiente, dos superan a la posición en la rejilla en alguna variable y una la iguala. GWA muestra asociación positiva con las \emph{epochs} hasta el umbral en 16 celdas y negativa en una, con mediana 0,40, y negativa con el \emph{accuracy} de test en 15 y positiva en una, con mediana $-0{,}29$. Una GWA alta en la ventana acompaña a un entrenamiento más lento y a un \emph{accuracy} final menor, en los cuatro conjuntos de datos y en las tres arquitecturas. La única celda de signo contrario es, en las dos variables, la ResNet-18 sobre Tiny-ImageNet con SGD. La \emph{gradient disparity} muestra asociación positiva con el \emph{accuracy} de test en 12 celdas y negativa en cuatro, con mediana 0,21. Las seis celdas de la convolucional y cinco de la red plana van en el mismo sentido. Las cuatro contrarias son dos de la red plana, sobre CIFAR-10 y sobre MNIST, y dos de la ResNet-18, sobre MNIST y sobre Tiny-ImageNet. La \emph{stiffness} muestra asociación negativa con el \emph{gap} de \emph{loss} en 15 celdas y positiva en una, con mediana $-0{,}38$. Una \emph{stiffness} alta acompaña a un \emph{gap} menor en las 13 celdas con asociación de la red plana y de la convolucional, y la ResNet-18 se parte. Son 15 celdas frente a las 15 de la posición en la rejilla, y la \emph{stiffness} iguala al \emph{learning rate} leído de otra forma, sin superarlo. Las otras tres métricas no superan a la posición en ninguna variable. GSNR empata con ella a siete celdas en velocidad y en \emph{accuracy} de test, y la escala de ruido a siete en velocidad; en los tres empates la posición queda delante por la mediana de $|D|$. En el resto de casos las tres quedan por debajo, con entre tres y 11 celdas. Su signo cambia además con el conjunto de datos o con la arquitectura. GSNR va con la velocidad en un sentido en CIFAR-10 y en el contrario en MNIST, y la escala de ruido se asocia con la velocidad sobre todo en la ResNet-18 y casi nada en la red plana.

La cuenta granulada deja la tabla casi vacía. Compara solo entrenamientos con el mismo \emph{learning rate}, es decir, las cinco \emph{seeds} de cada uno, y el \emph{learning rate} no puede mover la D. Ninguna de las 30 combinaciones muestra asociación en más de tres celdas, y las celdas que la muestran no coinciden de una variable a otra ni de una métrica a otra. Los recuentos de signo bajan a lo que una moneda da con frecuencia. Para las métricas de gradiente, el signo más frecuente reúne entre 10 y 16 celdas, con una excepción sin intervalo, la \emph{gradient confusion} con el \emph{accuracy} de test, 18 de 24. Las lecturas del predictor de referencia llegan a 18 y 19 con esa misma variable y con la velocidad, también sin intervalo fuera del cero en más de tres celdas. Las medianas del valor absoluto bajan a entre 0,10 y 0,33 en las seis métricas, y a 0,43 como mucho en el predictor de referencia, frente a los 0,18 a 0,71 de la cuenta por celda. La posición en la rejilla da cero en todas, porque dentro de un \emph{learning rate} no tiene nada que ordenar. Dentro de un \emph{learning rate}, ninguna métrica ordena las cinco \emph{seeds} mucho mejor que el azar, y eso vale también para el predictor de referencia. Lo que la cuenta por celda recoge es, en su mayor parte, que el \emph{learning rate} mueve a la vez la métrica y la eficiencia, según anticipó la Sección~\ref{sec:res-rango}.

H1 se sostiene como requisito previo, y solo por GWA. GWA muestra asociación con la velocidad en 16 celdas y con el \emph{accuracy} de test en 15, frente a las siete y siete de la posición en la rejilla, y con el mismo signo en los cuatro conjuntos de datos. Contando también las celdas del signo contrario, la ventaja se mantiene, 17 frente a 12 y 16 frente a 13. Las dos cuentas se dan porque la posición en la rejilla no cambia con la \emph{seed}. Su intervalo no es comparable con el de una métrica, y el requisito no puede depender de cómo se lea ese intervalo. La \emph{gradient disparity} supera también a la posición con el \emph{accuracy} de test, 12 frente a siete, pero su signo se da la vuelta en MNIST, con dos celdas negativas frente a una positiva, allí donde la variable apenas tiene recorrido. Con eso, las cinco hipótesis restantes tienen objeto. La asociación que sostiene H1 es, con todo, la que el \emph{learning rate} produce en las dos variables a la vez, y desaparece al fijarlo. Si esa asociación añade algo a lo que el \emph{loss} y el \emph{accuracy} de validación ya dicen en la misma ventana, sin coste, es la pregunta de H2. La D entre cada métrica y su referencia dentro de la celda ya avisa de que GWA y la \emph{stiffness} ordenan los entrenamientos en parte como el \emph{accuracy} de validación, con medianas de $-0{,}45$ y $-0{,}44$.

\section{Valor incremental sobre la referencia}
\label{sec:res-incremental}
% Fase C, escrita el 2026-09-04. Números de results/incremental.parquet y results/seleccion.parquet
% (2026-09-04), regenerables con `uv run python src/contrast.py`; figura `seleccion` en src/figures.py.

Esta sección responde a la hipótesis H2, la decisiva, con las tres lecturas de la Sección~\ref{sec:protocolo-objetivos}. La familia es la misma, la ventana del 5\,\% y los entrenamientos que aprendieron. Cada métrica se compara con la referencia nombrada de antemano para su variable, el \emph{accuracy} de validación para el \emph{accuracy} de test y el \emph{loss} de validación para el \emph{gap} de \emph{loss}. La primera lectura compara, celda a celda, el valor absoluto de la D de la métrica con el de su referencia. Cuenta en cuántas celdas la métrica gana y en cuántas pierde, y en cuántas de cada lado la diferencia deja fuera el cero por su intervalo \emph{jackknife}, calculado sobre los mismos entrenamientos. La segunda lee la D entre la métrica y su referencia dentro de la celda, $D_{\mathrm{ref}}$. Una métrica es redundante, es decir, la referencia con otro nombre, cuando el valor absoluto de esa D llega a 0,8, el umbral de la poda. La tercera es la lectura de selección de la Sección~\ref{sec:protocolo-seleccion}, la pérdida que la literatura de búsqueda de hiperparámetros llama \emph{regret}. En cada celda y \emph{seed} se elige el \emph{learning rate} que el predictor al 5\,\% señala como mejor, y se mide cuánto \emph{accuracy} de test se pierde frente al mejor \emph{learning rate} de esa \emph{seed}. La tabla da al lado, para cada predictor, en cuántas celdas esa pérdida es menor que la del \emph{accuracy} de validación, la columna «$<$ valid.». La posición en la rejilla no entra en esta lectura, porque elegir siempre el \emph{learning rate} más alto o el más bajo no es una predicción. La Tabla~\ref{tab:h2} recoge las tres lecturas para las dos variables que deciden, y la Figura~\ref{fig:seleccion} la tercera.

\begin{table}[ht]
\centering
\footnotesize
\setlength{\tabcolsep}{3.5pt}
\caption{H2 en la familia primaria. En negrita, la mejor métrica de gradiente en los recuentos crudos y en la selección.}
\label{tab:h2}
\begin{tabular}{lcccccccc}
\toprule
& \multicolumn{3}{c}{\textbf{\emph{Accuracy} de test}} & \multicolumn{3}{c}{\textbf{\emph{Gap} de \emph{loss}}} & \multicolumn{2}{c}{\textbf{Selección}} \\
\cmidrule(lr){2-4}\cmidrule(lr){5-7}\cmidrule(lr){8-9}
\textbf{Predictor} & gana/pierde & interv. & $|D_{\mathrm{ref}}|$ & gana/pierde & interv. & $|D_{\mathrm{ref}}|$ & pérdida & $<$ valid. \\
\midrule
Posición en la rejilla & 11/13 & 3/6 & 0,40 & 7/16 & 1/5 & 0,46 & & \\
\emph{Loss} de validación & 15/7 & 0/0 & 0,93 & ref. & ref. & ref. & 0,015 & 5 \\
\emph{Accuracy} de validación & ref. & ref. & ref. & 14/9 & 2/0 & 0,92 & 0,013 & ref. \\
TSE & 4/19 & 0/3 & 0,83 & 14/9 & 3/1 & 0,81 & 0,015 & 6 \\
\addlinespace
GNS & 6/18 & 0/8 & 0,26 & 3/21 & 0/12 & 0,25 & 0,073 & 0 \\
GSNR & 6/18 & 1/9 & 0,32 & 5/19 & 0/10 & 0,34 & 0,035 & 3 \\
\emph{Gradient disparity} & 10/14 & 3/8 & 0,36 & 5/18 & 0/7 & 0,37 & 0,097 & 1 \\
\addlinespace
\emph{Stiffness} & 5/19 & 0/8 & 0,44 & \textbf{7/17} & 2/6 & 0,44 & 0,110 & 1 \\
\emph{Gradient confusion} & 6/17 & 2/12 & 0,23 & 6/18 & 0/7 & 0,46 & 0,059 & 2 \\
GWA & \textbf{11/13} & 0/8 & 0,45 & 5/19 & 0/14 & 0,41 & \textbf{0,023} & \textbf{9} \\
\bottomrule
\end{tabular}
\end{table}

Sobre el \emph{accuracy} de test ninguna métrica gana a la validación. Con el recuento crudo la mejor es GWA, que gana en 11 celdas y pierde en 13, seguida de la \emph{gradient disparity} con 10 y 14; las otras cuatro ganan en cinco o seis. Con el intervalo el empate desaparece. GWA no gana con intervalo en ninguna celda y pierde en ocho. La disparidad es la única con tres celdas ganadas con intervalo, y pierde en ocho. Ninguna métrica es la validación con otro nombre. La mediana de $|D_{\mathrm{ref}}|$ va de 0,23 a 0,45, y solo dos celdas, las dos de GSNR, llegan a 0,8. Las métricas ordenan los entrenamientos de otra manera que la validación, y peor. Sin las seis celdas de MNIST, donde la variable apenas tiene recorrido, GWA queda en nueve celdas ganadas y nueve perdidas, ninguna con intervalo a favor y cinco en contra. La posición en la rejilla, que no hay que medir, hace lo mismo que la mejor métrica, 11 frente a 13. Dentro del predictor de referencia, las dos lecturas de la validación son intercambiables, con $|D_{\mathrm{ref}}|$ de 0,93 y las 24 celdas en 0,8 o más. TSE es redundante con el \emph{accuracy} de validación en 17 celdas y pierde contra él en 19. Los tres son una sola lectura.

Sobre el \emph{gap} de \emph{loss} la distancia es mayor. La referencia, el \emph{loss} de validación en la ventana, tiene una $|D|$ mediana de 0,54, y las seis métricas quedan entre 0,21 y 0,39. La \emph{stiffness} gana en siete celdas y pierde en 17, con dos ganadas y seis perdidas con intervalo. GWA gana en cinco y pierde en 19, con 14 perdidas con intervalo. El único predictor que supera a la referencia en el \emph{gap} es TSE, el \emph{loss} de entrenamiento acumulado, con 14 celdas frente a nueve y tres con intervalo frente a una. Los dos son gratuitos y los dos son lecturas de la velocidad, que a presupuesto fijo es en buena parte lo que el \emph{gap} mide, según la Sección~\ref{sec:variables}.

\begin{figure}[ht]
\centering
\includegraphics{seleccion}
\caption{Pérdida de \emph{accuracy} de test por selección al 5\,\%.}
\label{fig:seleccion}
\end{figure}

La lectura de selección, en la Figura~\ref{fig:seleccion}, pone la misma respuesta en puntos de \emph{accuracy}. Cada barra es la mediana sobre las 24 celdas y cada punto una celda. Elegir el \emph{learning rate} con el \emph{accuracy} de validación al 5\,\% pierde 0,013 frente al mejor de la celda; con el \emph{loss} de validación se pierde 0,015 y con TSE 0,015. Elegir al azar pierde 0,070. La mejor métrica es GWA, con 0,023, que pierde menos que la validación en nueve celdas repartidas entre los cuatro conjuntos. GSNR pierde 0,035 y la \emph{gradient confusion} 0,059. Las otras tres, con el signo que su artículo llama bueno, pierden más que elegir al azar, la escala de ruido 0,073, la disparidad 0,097 y la \emph{stiffness} 0,110. La cuenta con el signo invertido es una lectura a posteriori, sin regla previa, y se da como exploratoria. Con ella la disparidad baja a 0,023, la escala de ruido a 0,066 y la \emph{stiffness} a 0,075, y GWA, GSNR y la confusión suben a 0,084 o más. Para tres métricas el signo del artículo es el correcto en estas condiciones y para las otras tres es el contrario. Es lo mismo que H1 dejó ver en la disparidad, cuya D con el \emph{accuracy} de test es positiva.

Sobre la velocidad, donde H2 no decide, la referencia gana en todas partes. Ninguna métrica supera al \emph{accuracy} de validación al 5\,\% en más de cuatro celdas, y con intervalo pierden entre 13 y 19. El \emph{loss} de validación empata con él, 10 frente a 11, y TSE pierde en 21. Es lo que anuncia la Sección~\ref{sec:protocolo-eval}. La referencia es el primer tramo de la misma curva cuyo cruce se predice, y lleva ventaja por construcción.

H2 cae. En la ventana del 5\,\% ninguna métrica de gradiente supera a la lectura gratuita de la validación sobre el \emph{accuracy} de test ni sobre el \emph{gap}, ni con el intervalo de la diferencia ni eligiendo el \emph{learning rate} con ella. La que más se acerca, GWA, empata en el recuento crudo, pierde con intervalo y cuesta 0,023 de \emph{accuracy} frente a 0,013. Las métricas no son la validación con otro nombre, y aun así no la superan, porque lo que ordenan es el \emph{learning rate}, igual que la posición en la rejilla, que empata con la mejor de ellas. Si añaden algo encima de la validación, una vez fijada esta, no lo mide este diseño. La única cuenta libre del \emph{learning rate}, la granulada, queda vacía para todas, y también para la validación. El veredicto vale para esta ventana y esta rejilla. Qué pasa con ventanas posteriores lo mide la Sección~\ref{sec:res-temporal}, y lo que cuesta la instrumentación, la Sección~\ref{sec:res-coste}.

\section{Coste y capacidad predictiva}
\label{sec:res-coste}
% Fase C, escrita el 2026-09-04. Clases de coste de tab:coste-metricas; sobrecoste medido de
% desarrollo.tex §Cálculo de las métricas; capacidad predictiva de tab:h2.

La Tabla~\ref{tab:coste-metricas} separa los predictores en tres clases de coste. La clase gratuita reúne el \emph{loss} y el \emph{accuracy} de validación, TSE y la posición en la rejilla, que el entrenamiento produce sin calcular ningún gradiente. La clase del barrido compartido reúne cinco de las seis métricas de gradiente que quedan tras la poda. Las cinco calculan el gradiente de cada ejemplo del \emph{batch} de medición una sola vez, en un barrido común, y se separan después por la aritmética que hacen con él. Esa aritmética crece con el número de parámetros en la escala de ruido y GSNR, con el tamaño del \emph{batch} en GWA, y con el cuadrado del tamaño del \emph{batch} en la \emph{stiffness} y la \emph{gradient confusion}. La tercera clase es la de las submuestras, y solo tiene a la \emph{gradient disparity}, que sigue pagando sus cinco pasadas propias. Medido en reloj y sumado sobre los 960 entrenamientos, la instrumentación costó una cuarta parte del tiempo de entrenamiento. Por celda va del 0,7\,\% en la red convolucional sobre Tiny-ImageNet al 105\,\% en la red \emph{fully connected} sobre CIFAR-100 con SGD, es decir, hasta el doble del tiempo de entrenar, según mide la Sección~\ref{sec:des-metricas}. Ese reloj no separa una métrica de otra, porque el barrido es común.

Contra ese coste, la Tabla~\ref{tab:h2} no muestra ninguna métrica que supere a la clase gratuita. Dentro de la clase del barrido, la mejor sobre el \emph{accuracy} de test es GWA, la de aritmética más ligera. Empata con la posición en la rejilla y pierde con intervalo contra la validación. Sobre el \emph{gap} la mejor es la \emph{stiffness}, de coste cuadrático, con siete celdas ganadas frente a 17 perdidas. La disparidad, la única que paga pasadas propias, casi iguala a GWA en el recuento crudo del \emph{accuracy} de test y pierde en todo lo demás. Pagar más aritmética no compra orden. Eso responde a OE2. En la ventana del 5\,\% ninguna clase de coste compra más capacidad predictiva que la clase gratuita, y el sobrecoste, hasta el doble del tiempo de entrenamiento en la red plana, no está justificado como predicción. Lo que sí justifica es medir, que es lo que este trabajo hace con él, porque la serie completa por \emph{epoch} es la que permite responder a las hipótesis restantes.

\section{Qué familia predice mejor}
\label{sec:res-familia}
% Fase D, escrita el 2026-09-04. Números de results/ranking.parquet y ranking_grupos.parquet
% (2026-09-04), regenerables con `uv run python src/contrast.py`; las tres parejas de la poda las
% imprime `uv run python src/efficiency.py`. Sin figura: una comparación de niveles va a tabla.

Esta sección responde a la hipótesis H3 con la regla de la Sección~\ref{sec:protocolo-objetivos}. Para cada variable de la familia primaria, las seis métricas de gradiente se ordenan en un \emph{ranking} por el número de celdas en las que muestran asociación con su signo mayoritario, y a igualdad por la mediana del valor absoluto de D. No se calcula nada nuevo, porque las tres cifras salen de la Tabla~\ref{tab:signos-h1}, y la familia de cada métrica solo etiqueta el orden. La Tabla~\ref{tab:h3} da los tres órdenes, con la posición en la rejilla intercalada en el puesto que le tocaría y sin número. Contar solo las celdas del signo mayoritario penaliza a la métrica que cambia de sentido de una celda a otra, y debe hacerlo. Un predictor que ordena en un sentido en unas celdas y en el contrario en otras no sirve sin saber antes en qué celda se está. La escala de ruido, por ejemplo, muestra asociación negativa con el \emph{accuracy} de test en seis celdas y positiva en tres, pero su D es positiva en 14 celdas de 24, así que su signo mayoritario es el positivo y puntúa tres. La regla de tomar el signo mayoritario sobre las 24 celdas, en lugar de solo sobre las que muestran asociación, se fijó con la Tabla~\ref{tab:signos-h1} a la vista. La otra lectura subiría a la escala de ruido con el \emph{accuracy} de test del sexto al quinto puesto y dejaría igual los tres primeros.

\begin{table}[ht]
\centering
\footnotesize
\setlength{\tabcolsep}{4pt}
\caption{Orden de las métricas de gradiente por variable al 5\,\%, con las celdas a favor y en contra del signo mayoritario.}
\label{tab:h3}
\begin{tabular}{clcccc}
\toprule
\textbf{Puesto} & \textbf{Predictor} & \textbf{Familia} & \textbf{Signo} & \textbf{Celdas} & \textbf{Mediana de $|D|$} \\
\midrule
\multicolumn{6}{l}{\textbf{\emph{Epochs} hasta el umbral}} \\
1 & GWA & alineación & $+$ & 16/1 & 0,40 \\
2 & \emph{Stiffness} & alineación & $+$ & 9/5 & 0,35 \\
3 & \emph{Gradient disparity} & variabilidad & $-$ & 9/4 & 0,28 \\
  & Posición en la rejilla & & $-$ & 7/5 & 0,33 \\
4 & GSNR & variabilidad & $-$ & 7/6 & 0,33 \\
5 & GNS & variabilidad & $+$ & 7/3 & 0,29 \\
6 & \emph{Gradient confusion} & alineación & $+$ & 5/4 & 0,24 \\
\addlinespace
\multicolumn{6}{l}{\textbf{\emph{Accuracy} de test}} \\
1 & GWA & alineación & $-$ & 15/1 & 0,29 \\
2 & \emph{Gradient disparity} & variabilidad & $+$ & 12/4 & 0,34 \\
  & Posición en la rejilla & & $-$ & 7/6 & 0,30 \\
3 & GSNR & variabilidad & $+$ & 7/4 & 0,24 \\
4 & \emph{Stiffness} & alineación & $+$ & 5/6 & 0,24 \\
5 & \emph{Gradient confusion} & alineación & $-$ & 5/5 & 0,18 \\
6 & GNS & variabilidad & $+$ & 3/6 & 0,22 \\
\addlinespace
\multicolumn{6}{l}{\textbf{\emph{Gap} de \emph{loss}}} \\
  & Posición en la rejilla & & $+$ & 15/0 & 0,46 \\
1 & \emph{Stiffness} & alineación & $-$ & 15/1 & 0,39 \\
2 & \emph{Gradient disparity} & variabilidad & $-$ & 11/4 & 0,36 \\
3 & GSNR & variabilidad & $-$ & 11/2 & 0,32 \\
4 & \emph{Gradient confusion} & alineación & $+$ & 9/5 & 0,31 \\
5 & GNS & variabilidad & $-$ & 6/5 & 0,25 \\
6 & GWA & alineación & $-$ & 4/6 & 0,21 \\
\bottomrule
\end{tabular}
\end{table}

El primer puesto es de la familia de alineación en las tres variables. El segundo es de la de variabilidad con el \emph{accuracy} de test y con el \emph{gap}, y en velocidad lo ocupa la \emph{stiffness}, que es de alineación. GWA encabeza la velocidad y el \emph{accuracy} de test, con 16 y 15 celdas, y la \emph{stiffness} el \emph{gap}, con 15. Detrás va la \emph{gradient disparity}, con nueve, 12 y 11 celdas, empatada a nueve con la \emph{stiffness} en velocidad y por detrás de ella por la mediana. Las otras dos métricas de alineación no acompañan a las dos primeras. La \emph{gradient confusion} es la última en velocidad y la penúltima en el \emph{accuracy} de test. GWA, la mejor en dos variables, es la última en el \emph{gap}, con cuatro celdas a favor y seis en contra. La posición en la rejilla queda por debajo de tres métricas en velocidad y de dos en el \emph{accuracy} de test, y en el \emph{gap} por encima de las seis, con 15 celdas y una mediana de 0,46 frente a 0,39 de la \emph{stiffness}. Ninguna métrica se acerca a las lecturas del predictor de referencia, que con la misma regla puntúan entre 16 y 24 celdas.

El orden cambia de bando con la arquitectura y con el conjunto de datos. Repetida la cuenta dentro de cada arquitectura, sobre ocho celdas, en la red \emph{fully connected} gana la alineación en las tres variables. GWA gana en velocidad y en \emph{accuracy} de test y la \emph{stiffness} en el \emph{gap}, con siete de las ocho celdas y ninguna en contra en los tres casos. En la convolucional gana la variabilidad en velocidad y en \emph{accuracy} de test, GSNR y la disparidad, y la \emph{stiffness} en el \emph{gap}. En la ResNet-18 ganan GSNR en velocidad, la escala de ruido en \emph{accuracy} de test y GWA en el \emph{gap}. Esa GWA es además una inversión de signo, porque con el \emph{gap} muestra asociación negativa en cuatro celdas de la convolucional, positiva en seis de la ResNet-18 y ninguna en la red plana. Por conjunto de datos, sobre seis celdas, en MNIST y en CIFAR-10 gana la alineación en las tres variables, GWA y la \emph{stiffness}. En CIFAR-100 gana la variabilidad en las tres, la escala de ruido, la disparidad y GSNR. Tiny-ImageNet se parte entre la \emph{stiffness}, la disparidad y la \emph{gradient confusion}. Lo que la Tabla~\ref{tab:h3} ordena sobre las 24 celdas es en buena parte lo que pasa en la red \emph{fully connected} y en los dos conjuntos más fáciles, donde la alineación separa los \emph{learning rates} con más claridad que en el resto.

La poda es también un resultado. De las ocho métricas medidas, la m-coherencia y la varianza normalizada son la escala de ruido con otro nombre o con otro estimador, y la \emph{gradient disparity} es una medida de variabilidad, según deriva la Sección~\ref{sec:protocolo-objetivos}. Las tres identidades se miden aquí con el estadístico del contraste, la D entre las dos columnas dentro de cada celda al 5\,\% sobre los entrenamientos que aprendieron. La escala de ruido y la m-coherencia, de familias distintas en sus artículos, dan $|D| = 1$ en las 24 celdas, la identidad exacta de la Ecuación~\ref{eq:gns-mcoh}. La varianza normalizada ordena los entrenamientos como la escala de ruido con una mediana de 0,84 y 18 celdas en 0,8 o más, el umbral con el que el protocolo declara redundante una métrica. Contra GSNR, con la que el primer criterio la comparó, se queda en 0,59 y en ninguna celda llega a 0,8. La \emph{gradient disparity}, que su artículo define como el desacuerdo entre los gradientes de dos submuestras, ordena como la raíz de la traza de la covarianza del ruido con una mediana de 0,93 y 23 celdas en 0,8 o más. Por eso el contraste la lee como variabilidad. Dos de las ocho métricas que la literatura presenta por separado no aportan sobre esta matriz una ordenación propia, una tercera cambia de familia, y la comparación entre familias queda en tres contra tres.

H3 cae en la forma en que se enunció. Ninguna familia ocupa las primeras posiciones. La alineación pone la mejor métrica de cada variable y también la peor en dos de las tres, y la variabilidad ocupa el segundo puesto en las tres. El orden tampoco se mantiene entre arquitecturas ni entre conjuntos, porque la alineación gana las tres variables en la red \emph{fully connected} y en MNIST y CIFAR-10, y la variabilidad dos en la convolucional y las tres en CIFAR-100. Lo que se puede sostener sobre la alineación de gradientes es más estrecho. La métrica que más celdas ordena en cada variable es una de alineación, GWA o la \emph{stiffness}, y las dos siguen por debajo de la validación gratuita según decide la Sección~\ref{sec:res-incremental}. Con eso se sostiene una afirmación sobre dos métricas; la familia entera queda sin respaldo.

\section{Suficiencia de la ventana temprana}
\label{sec:res-temporal}
% Fase E, escrita el 2026-09-04. Números de results/ventanas_recuento.parquet, ventanas.parquet e
% incremental.parquet (2026-09-04), regenerables con `uv run python src/contrast.py`; figura
% `ventanas` en src/figures.py.

Esta sección responde a la hipótesis H4, que pregunta si medir en la ventana del 5\,\% predice tan bien como medir más tarde. La cuenta es la de la Sección~\ref{sec:protocolo-objetivos}. Para las variables que se leen al terminar el entrenamiento, el \emph{accuracy} de test y el \emph{gap} de \emph{loss}, se compara la D de la ventana del 5\,\% con la D de la ventana del 50\,\%, celda a celda. Las dos D salen de los mismos entrenamientos, y su diferencia lleva el mismo intervalo \emph{jackknife} pareado que la Sección~\ref{sec:res-incremental}. Se cuentan las celdas en las que $|D|$ crece al esperar y las celdas en las que decrece, y cuántas de cada lado lo hacen con el intervalo fuera del cero. Para la velocidad solo sirven las dos primeras ventanas, y las dos se leen por hitos, entre los entrenamientos que aún no habían cruzado al cerrar cada ventana, según fija la Sección~\ref{sec:protocolo-hitos}. Esa comparación es entre las \emph{epochs} 1 y 2 en MNIST y entre la 2 y la 4 en los otros tres conjuntos. Como los entrenamientos en riesgo no son los mismos en las dos ventanas, ahí el intervalo suma las dos varianzas y es más ancho. Al cerrar la segunda ventana quedan en riesgo entre 8 y 40 entrenamientos por celda. La cuenta de velocidad es sobre 23 celdas. En Tiny-ImageNet con la red \emph{fully connected} bajo Adam ninguno de los diez entrenamientos en riesgo cruza después. Sin ningún par comparable, esa celda no entra. La Tabla~\ref{tab:h4} da, para cada predictor y variable, las celdas que suben y las que bajan, cuántas de cada lado tienen el intervalo fuera del cero, y la mediana de $|D|$ en la ventana temprana y en la tardía. La Figura~\ref{fig:ventanas} dibuja el cambio de $|D|$ de cada celda, resumido en cajas como en la Figura~\ref{fig:signos-h1}.

\begin{table}[ht]
\centering
\caption{H4 en la familia primaria, de la ventana temprana a la tardía.}
\label{tab:h4}
\scriptsize
\setlength{\tabcolsep}{2.5pt}
\begin{tabular}{lccccccccc}
\toprule
& \multicolumn{3}{c}{\textbf{\emph{Accuracy} de test}} & \multicolumn{3}{c}{\textbf{\emph{Gap} de \emph{loss}}} & \multicolumn{3}{c}{\textbf{\emph{Epochs} hasta el umbral}} \\
\cmidrule(lr){2-4}\cmidrule(lr){5-7}\cmidrule(lr){8-10}
\textbf{Predictor} & sube/baja & interv. & $|D|$ & sube/baja & interv. & $|D|$ & sube/baja & interv. & $|D|$ \\
\midrule
\emph{Loss} de validación & 14/10 & 8/6 & 0,39$\to$0,52 & 6/18 & 4/4 & 0,54$\to$0,42 & 14/9 & 1/0 & 0,71$\to$0,74 \\
\emph{Accuracy} de validación & 22/2 & 16/0 & 0,40$\to$0,70 & 4/20 & 1/8 & 0,53$\to$0,33 & 11/12 & 1/0 & 0,71$\to$0,69 \\
TSE & 18/5 & 13/1 & 0,33$\to$0,50 & 10/14 & 6/6 & 0,51$\to$0,61 & 9/14 & 0/0 & 0,64$\to$0,59 \\
\addlinespace
GNS & 7/16 & 0/0 & 0,22$\to$0,21 & 10/13 & 0/1 & 0,25$\to$0,13 & 6/17 & 0/1 & 0,29$\to$0,19 \\
GSNR & 13/11 & 2/1 & 0,24$\to$0,30 & 9/15 & 4/1 & 0,32$\to$0,25 & 9/14 & 0/1 & 0,33$\to$0,27 \\
\emph{Gradient disparity} & 11/13 & 6/4 & 0,34$\to$0,35 & 10/13 & 4/4 & 0,36$\to$0,30 & 19/4 & 1/0 & 0,28$\to$0,44 \\
\addlinespace
\emph{Stiffness} & 15/9 & 4/2 & 0,24$\to$0,27 & 14/9 & 5/3 & 0,39$\to$0,51 & 10/13 & 0/0 & 0,35$\to$0,28 \\
\emph{Gradient confusion} & 16/8 & 5/3 & 0,18$\to$0,29 & 12/12 & 4/1 & 0,31$\to$0,41 & 17/6 & 2/0 & 0,24$\to$0,30 \\
GWA & 12/12 & 2/1 & 0,29$\to$0,33 & 11/12 & 1/3 & 0,21$\to$0,16 & 11/12 & 1/0 & 0,40$\to$0,41 \\
\bottomrule
\end{tabular}
\end{table}

\begin{figure}[ht]
\centering
\includegraphics{ventanas}
\caption{Cambio de $|D|$ por celda al esperar.}
\label{fig:ventanas}
\end{figure}

Sobre el \emph{accuracy} de test, las seis métricas de gradiente no ganan nada esperando. Con el intervalo, la que más crece es la \emph{gradient disparity}, en seis celdas, y ninguna decrece en más de cuatro. Los recuentos crudos se reparten cerca de la mitad, entre 7 y 16 celdas que crecen de 24. Las medianas de $|D|$ van de 0,18 a 0,34 al 5\,\% y de 0,21 a 0,35 al 50\,\%. La que más sube es la \emph{gradient confusion}, de 0,18 a 0,29, con 16 celdas en crudo y cinco con intervalo. La \emph{stiffness} crece en 15 celdas en crudo y en cuatro con intervalo. Lo que sí crece es el predictor de referencia. El \emph{accuracy} de validación pasa de 0,40 a 0,70, con 22 celdas que crecen y 16 con intervalo, y TSE de 0,33 a 0,50, con 18 y 13. A mitad del entrenamiento la curva de validación ya es casi el resultado que se quiere predecir, y por eso su D crece. La Figura~\ref{fig:ventanas} lo enseña de un vistazo. En el panel del \emph{accuracy} de test, las cajas del \emph{accuracy} de validación y de TSE quedan a la derecha del cero, y las cajas de las seis métricas quedan centradas sobre él. La comparación de H2 repetida al 50\,\% dice lo mismo. En esa ventana ninguna métrica supera a la validación en más de dos celdas, y GWA pierde en 23, frente a las 13 del 5\,\%.

Sobre el \emph{gap} de \emph{loss} ocurre lo contrario con el predictor de referencia. El \emph{accuracy} de validación pierde $|D|$ al esperar, de 0,53 a 0,33, con 20 celdas que decrecen y ocho con intervalo, y el \emph{loss} de validación baja de 0,54 a 0,42. A presupuesto fijo el \emph{gap} es en buena parte velocidad leída al revés, según la Sección~\ref{sec:variables}, y la lectura temprana de la validación es la que mejor la recoge. Las métricas se reparten otra vez cerca de la mitad. Solo la \emph{stiffness} y la \emph{gradient confusion} crecen algo, de 0,39 a 0,51 y de 0,31 a 0,41, con cinco y cuatro celdas con intervalo. TSE es el único predictor que supera con claridad al \emph{loss} de validación al 50\,\%, con 16 celdas frente a ocho y nueve con intervalo, y sigue sin costar nada. El \emph{gap} de \emph{accuracy}, la lectura de apoyo, no cambia el cuadro, porque ninguna métrica crece con intervalo en más de cinco celdas.

Sobre la velocidad, entre la primera y la segunda ventana, casi nada se mueve con intervalo. Cada predictor cambia con intervalo en dos celdas como mucho, y las poblaciones en riesgo bajan hasta ocho entrenamientos en la segunda ventana. Los dos recuentos crudos que destacan son la \emph{gradient disparity}, que crece en 19 celdas y pasa de 0,28 a 0,44, y la \emph{gradient confusion}, que crece en 17. El resto se queda donde estaba. Las dos lecturas de la validación se quedan entre 0,69 y 0,74 en las dos ventanas, y TSE baja de 0,64 a 0,59.

H4 se sostiene para las métricas de gradiente, y por la razón menos favorable. Medir al 5\,\% predice tan bien como medir al 50\,\% porque su $|D|$ no crece con la ventana en la mayoría de las celdas con intervalo. Sobre el \emph{accuracy} de test, cuatro de las seis métricas se quedan planas, con medianas entre 0,2 y 0,35 en las dos ventanas. La \emph{stiffness} y la \emph{gradient confusion} crecen algo con el \emph{gap}, hasta 0,51 y 0,41, sin pasar de cinco celdas con intervalo. Ninguna se acerca a la validación, que sube hasta 0,70 con el \emph{accuracy} de test sin costar nada. La ventana mínima defendible es, por tanto, la primera, la \emph{epoch} 1 en MNIST y la 2 en los demás conjuntos. Lo es porque cuatro métricas no maduran y dos maduran poco, y esperar no cambia lo que dicen. Si bastaría medir antes no lo puede decir esta matriz, porque no tiene ninguna lectura anterior.

\section{Robustez entre optimizadores}
\label{sec:res-robustez}
% Fase F, escrita el 2026-09-04. Números de results/optimizadores.parquet (2026-09-04),
% regenerable con `uv run python src/contrast.py`. Sin figura: las inversiones van de 0 a 2.

Esta sección responde a la hipótesis H5 con la regla de la Sección~\ref{sec:protocolo-objetivos}. Las 24 celdas forman 12 pares. Las dos celdas de un par comparten conjunto de datos y arquitectura y solo difieren en el optimizador. El diseño las hizo comparables a propósito, con las mismas cinco \emph{seeds}, un umbral compartido y una rejilla de la misma forma, la de Adam la de SGD dividida entre 10. Para cada predictor y variable de la familia primaria se hacen tres recuentos sobre los 12 pares. El primero cuenta los pares que conservan el signo de la D y el segundo los que lo invierten. En los dos solo entran los pares en los que las dos D dejan fuera el cero por su intervalo, porque dos signos que coinciden sin intervalo coinciden por azar la mitad de las veces. Un par en el que alguna de las dos D incluye el cero no cuenta a ningún lado. El tercero cuenta los pares en los que la diferencia entre la D con SGD y la D con Adam deja fuera el cero. Como las dos D salen de entrenamientos distintos, ese intervalo suma las dos varianzas. Este recuento recoge los cambios de tamaño que no llegan a cambiar el signo. La Tabla~\ref{tab:h5} da los tres recuentos, como ac./inv. y dif., y la mediana del valor absoluto de D en cada brazo, como SGD/Adam.

\begin{table}[ht]
\centering
\footnotesize
\setlength{\tabcolsep}{2.5pt}
\caption{H5 sobre los 12 pares de celdas, $|D|$ con SGD y con Adam. En negrita, las inversiones.}
\label{tab:h5}
\begin{tabular}{lccccccccc}
\toprule
& \multicolumn{3}{c}{\textbf{\emph{Epochs} hasta el umbral}} & \multicolumn{3}{c}{\textbf{\emph{Accuracy} de test}} & \multicolumn{3}{c}{\textbf{\emph{Gap} de \emph{loss}}} \\
\cmidrule(lr){2-4}\cmidrule(lr){5-7}\cmidrule(lr){8-10}
\textbf{Predictor} & ac./inv. & dif. & SGD/Adam & ac./inv. & dif. & SGD/Adam & ac./inv. & dif. & SGD/Adam \\
\midrule
Posición en la rejilla & 1/\textbf{1} & 6 & 0,30/0,36 & 2/\textbf{2} & 6 & 0,42/0,18 & 4/0 & 3 & 0,51/0,40 \\
\addlinespace
\emph{Loss} de validación & 12/0 & 3 & 0,69/0,78 & 7/0 & 5 & 0,38/0,39 & 10/0 & 5 & 0,61/0,45 \\
\emph{Accuracy} de validación & 12/0 & 3 & 0,63/0,77 & 6/0 & 4 & 0,42/0,38 & 10/0 & 7 & 0,62/0,46 \\
TSE & 11/0 & 3 & 0,58/0,73 & 6/0 & 5 & 0,28/0,36 & 11/0 & 6 & 0,56/0,51 \\
\addlinespace
GNS & 4/0 & 0 & 0,27/0,30 & 3/0 & 0 & 0,23/0,22 & 1/\textbf{1} & 3 & 0,26/0,16 \\
GSNR & 5/0 & 1 & 0,27/0,35 & 4/0 & 0 & 0,20/0,28 & 4/0 & 5 & 0,25/0,33 \\
\emph{Gradient disparity} & 3/\textbf{1} & 4 & 0,48/0,23 & 5/\textbf{2} & 4 & 0,34/0,30 & 4/\textbf{2} & 5 & 0,37/0,31 \\
\addlinespace
\emph{Stiffness} & 3/\textbf{1} & 3 & 0,43/0,29 & 2/\textbf{1} & 4 & 0,25/0,15 & 5/0 & 2 & 0,21/0,43 \\
\emph{Gradient confusion} & 1/0 & 1 & 0,21/0,27 & 2/0 & 2 & 0,18/0,18 & 4/0 & 2 & 0,35/0,29 \\
GWA & 6/0 & 2 & 0,37/0,45 & 6/0 & 4 & 0,27/0,29 & 3/0 & 3 & 0,19/0,30 \\
\bottomrule
\end{tabular}
\end{table}

Ninguna métrica invierte el signo en más de dos pares de doce. La \emph{gradient disparity} invierte uno con la velocidad y dos con el \emph{accuracy} de test y con el \emph{gap}. La \emph{stiffness} invierte uno con la velocidad y uno con el \emph{accuracy} de test, y la escala de ruido uno con el \emph{gap}. GSNR, la \emph{gradient confusion} y GWA no invierten ninguno. Las inversiones de una misma métrica se concentran en pocos pares. Las de la disparidad con la velocidad y con el \emph{accuracy} de test caen en el mismo par, la red \emph{fully connected} sobre CIFAR-10. La segunda del \emph{accuracy} de test es la ResNet-18 sobre MNIST, y las dos del \emph{gap} son las redes convolucionales sobre CIFAR-100 y sobre Tiny-ImageNet. Las dos de la \emph{stiffness} caen en un solo par, la ResNet-18 sobre CIFAR-100. En esos pares la métrica cambia de sentido con las dos variables a la vez. Eso es lo que hace una métrica que lee el \emph{learning rate} cuando el mejor \emph{learning rate} de la rejilla queda a un lado con SGD y al otro con Adam. La posición en la rejilla lo muestra sola, con sus tres inversiones en MNIST. Los pares que acuerdan con intervalo son pocos para las métricas, entre uno y seis de doce. GWA es la que más acuerda en velocidad y en \emph{accuracy} de test, con seis pares en las dos, y la \emph{stiffness} en el \emph{gap}, con cinco. En el resto de pares una de las dos D incluye el cero y el par no cuenta. Las lecturas del predictor de referencia acuerdan en 11 o 12 pares con la velocidad, en seis o siete con el \emph{accuracy} de test y en 10 u 11 con el \emph{gap}, sin ninguna inversión.

Lo que sí cambia entre optimizadores es el tamaño de la asociación. La \emph{gradient disparity} pasa de una mediana de $|D|$ de 0,48 con SGD a 0,23 con Adam en velocidad, con cuatro pares en los que la diferencia deja fuera el cero. La \emph{stiffness} sube de 0,21 a 0,43 en el \emph{gap}. La posición en la rejilla baja de 0,42 a 0,18 con el \emph{accuracy} de test, con seis diferencias fuera del cero. Las lecturas del predictor de referencia son más fuertes con Adam en velocidad, entre 0,73 y 0,78 frente a entre 0,58 y 0,69, y más fuertes con SGD en el \emph{gap}, entre 0,56 y 0,62 frente a entre 0,45 y 0,51. Esas diferencias se leen primero como diferencias de población y solo después como efecto del optimizador. Los 41 entrenamientos divergidos son todos de SGD, y Adam no divergió en ninguno de sus 480. La norma del umbral de la Sección~\ref{sec:calibracion} se calculó sobre los dos optimizadores juntos, y la censura de la velocidad no es la misma en los dos brazos de un par.

H5 no se refuta y no queda demostrada. Con 12 pares, y entre uno y seis pares con las dos D fuera del cero por métrica, ninguna asociación segura cambia de signo en más de dos pares. Las inversiones que hay llegan de dos en dos, con las dos variables de un mismo par, con la forma de una métrica que lee el \emph{learning rate}. Este diseño no puede afirmar invariancia, porque 12 pares solo detectan inversiones groseras, y no encontrarlas no prueba que no las haya. Lo que sí se puede decir es que, allí donde una métrica ordena con seguridad bajo los dos optimizadores, ordena en el mismo sentido, y que el tamaño de esa ordenación cambia con el optimizador más de lo que cambia su signo.

\section{Concordancia con la literatura}
\label{sec:res-concordancia}
% Fase G, escrita el 2026-09-04 (H6, obj. \ref{obj:concordancia}). Números de results/signos.parquet
% (2026-09-04), regenerable con `uv run python src/contrast.py`; las predicciones son contrast.PREDICTED,
% la tabla de signos del vault por variable. Sin figura: son recuentos. La interpretación vive en Conclusiones.

Esta sección responde a la hipótesis H6 con la regla de la Sección~\ref{sec:protocolo-objetivos}. El artículo que propuso cada métrica predice el sentido de su relación con alguna variable, y la Sección~\ref{sec:hipotesis} fija esas predicciones. Para cada una se cuentan, sobre las 24 celdas de la familia primaria, las celdas con asociación del signo predicho y las celdas con asociación del signo contrario. El criterio de celda es el mismo intervalo de dos colas de las secciones anteriores, sin bajar la barra, y una celda significa lo mismo en todo el capítulo. No se calcula nada nuevo, porque las cifras son las de la Tabla~\ref{tab:signos-h1} partidas por el signo que cada artículo afirma. Entran siete predicciones de seis métricas. La \emph{stiffness} intra-clase y la \emph{gradient confusion} tienen predicción sobre la velocidad. La escala de ruido la tiene sobre la velocidad, donde coinciden la predicción de McCandlish et al. y la de la m-coherencia leída con el signo cambiado por la Ecuación~\ref{eq:gns-mcoh}, y sobre el \emph{gap}, donde la única predicción es la de la m-coherencia. GSNR la tiene sobre el \emph{gap}. La \emph{gradient disparity} la tiene sobre el \emph{accuracy} de test, porque su artículo la liga al error de test. GWA la tiene sobre el \emph{accuracy} de test, con el signo convertido de la Sección~\ref{sec:def-alineacion}. Las cinco predicciones que esta memoria extrapola de un artículo a otra variable ocupan la segunda mitad de la Tabla~\ref{tab:h6} y quedan fuera del veredicto. La tabla repite la cuenta dentro de cada arquitectura, porque la Sección~\ref{sec:res-familia} mostró que el signo cambia con ella. En negrita van las cuentas en contra que superan a las cuentas a favor.

\begin{table}[ht]
\centering
\footnotesize
\setlength{\tabcolsep}{3.5pt}
\caption{H6 al 5\,\%, sobre las 24 celdas y dentro de cada arquitectura. En negrita, las cuentas en contra que superan a las de a favor.}
\label{tab:h6}
\begin{tabular}{llllcccc}
\toprule
\textbf{Predictor} & \textbf{Variable} & \textbf{Base} & \textbf{Signo} & \textbf{Celdas} & \textbf{FC} & \textbf{CNN} & \textbf{ResNet-18} \\
\midrule
\multicolumn{8}{l}{\textbf{Predicciones que el artículo enuncia}} \\
\emph{Stiffness} & Velocidad & Fort et al.~\cite{fort2019stiffness} & $-$ & 5/\textbf{9} & 1/2 & 0/4 & 4/3 \\
\emph{Gradient confusion} & Velocidad & Sankararaman et al.~\cite{sankararaman2020impact} & $+$ & 5/4 & 4/0 & 0/3 & 1/1 \\
GNS & Velocidad & McCandlish et al.~\cite{mccandlish2018empirical} & $+$ & 7/3 & 0/1 & 2/0 & 5/2 \\
 & & y Chatterjee~\cite{chatterjee2020coherent} & & & & & \\
GNS & \emph{Gap} & Chatterjee y Zielinski~\cite{chatterjee2020making} & $+$ & 5/\textbf{6} & 0/2 & 1/3 & 4/1 \\
GSNR & \emph{Gap} & Liu et al.~\cite{liu2020understanding} & $-$ & 11/2 & 4/2 & 3/0 & 4/0 \\
\emph{Gradient disparity} & \emph{Accuracy} & Forouzesh y Thiran~\cite{forouzesh2021disparity} & $-$ & 4/\textbf{12} & 2/5 & 0/6 & 2/1 \\
GWA & \emph{Accuracy} & Hölzl et al.~\cite{hoelzl2025gradient} & $-$ & 15/1 & 7/0 & 3/0 & 5/1 \\
\addlinespace
\multicolumn{8}{l}{\textbf{Predicciones que esta memoria extrapola, sin contrastar}} \\
\emph{Gradient disparity} & Velocidad & del error de test & $+$ & 4/\textbf{9} & 2/5 & 0/4 & 2/0 \\
GSNR & Velocidad & del \emph{gap} & $-$ & 7/6 & 1/0 & 0/5 & 6/1 \\
GSNR & \emph{Accuracy} & del \emph{gap} & $+$ & 7/4 & 3/0 & 0/4 & 4/0 \\
GWA & Velocidad & del \emph{accuracy} de test & $+$ & 16/1 & 7/0 & 4/0 & 5/1 \\
GWA & \emph{Gap} & del \emph{accuracy} de test & $+$ & 6/4 & 0/0 & 0/4 & 6/0 \\
\bottomrule
\end{tabular}
\end{table}

Dos de las tres predicciones sobre el desenlace final se cumplen. GWA con el \emph{accuracy} de test va en el sentido de su artículo en 15 celdas y en contra en una, la ResNet-18 sobre Tiny-ImageNet con SGD. Lo hace en las tres arquitecturas, con siete de las ocho celdas de la red \emph{fully connected}. GSNR con el \emph{gap} va a favor en 11 celdas y en contra en dos, las dos de la red \emph{fully connected} sobre CIFAR-10, también en las tres arquitecturas. La tercera va en contra. El artículo de la \emph{gradient disparity} la liga a un error de test mayor. Aquí una disparidad alta acompaña a un \emph{accuracy} de test mayor en 12 celdas, las seis de la convolucional, cinco de la red \emph{fully connected} y una de la ResNet-18, frente a cuatro en el sentido del artículo. Las tres predicciones sobre la velocidad apenas reciben apoyo. La escala de ruido va a favor en siete celdas y en contra en tres, las tres en MNIST, y la sostiene la ResNet-18 con cinco de las siete. La \emph{gradient confusion} da cinco celdas a favor y cuatro en contra, y las dos cuentas no se mezclan, porque las cuatro de la red \emph{fully connected} van con su artículo y las tres de la convolucional en contra. La \emph{stiffness} va en contra, con nueve celdas frente a cinco. La convolucional la contradice en sus cuatro celdas con asociación, y la ResNet-18 la apoya en cuatro y la contradice en tres. La predicción de la m-coherencia sobre el \emph{gap}, leída sobre la escala de ruido, da cinco celdas a favor y seis en contra.

Donde el signo contradice al artículo, la Sección~\ref{sec:protocolo-objetivos} pide leer primero la composición de dos curvas sobre el \emph{learning rate}, y las dos discrepancias tienen esa forma. El artículo de la \emph{stiffness} fijó su signo a lo largo del tiempo de entrenamiento, con la \emph{stiffness} cayendo hacia cero al entrar en el sobreajuste. Aquí, dentro de la celda, la \emph{stiffness} del 5\,\% ordena los entrenamientos al revés que el \emph{accuracy} de validación de esa misma ventana, con una D de mediana $-0{,}44$ según la Sección~\ref{sec:res-por-condicion}. El entrenamiento que más ha avanzado en la ventana tiene los gradientes de una misma clase menos alineados, y es el que cruza el umbral antes. Una \emph{stiffness} alta marca, por tanto, a un entrenamiento retrasado por su \emph{learning rate}, y esa cuenta no pone a prueba el mecanismo del artículo. La \emph{gradient disparity} es el mismo caso con el signo contrario. Su artículo la vio crecer con el sobreajuste y con el ruido de etiquetas. Aquí ordena los entrenamientos como el \emph{accuracy} de validación de la ventana, con una D de mediana 0,26 y positiva en 16 celdas, porque las dos dibujan un pico a lo largo del \emph{learning rate}. Una disparidad alta marca al entrenamiento que más ha avanzado en la ventana, y el sobreajuste que su artículo describe es un eje que la celda no recorre. La \emph{gradient confusion}, cuyo artículo fijó el signo a lo largo de la anchura y la profundidad de la red, cambia de signo entre la red plana y la convolucional, con cuentas de cuatro y de tres celdas. Del lado de los acuerdos vale el aviso contrario. GWA coincide con su artículo casi sin excepción, pero ordena los entrenamientos en parte como el \emph{accuracy} de validación, con una D de mediana $-0{,}45$, y peor que él, según midió la Sección~\ref{sec:res-incremental}. La razón está en la Sección~\ref{sec:def-alineacion}. GWA es un margen normalizado que el paso hacia delante ya calcula casi por completo. Su acuerdo confirma al artículo en lo que el artículo mide, sin añadir una prueba independiente sobre los gradientes, porque no supera a la validación.

De las extrapoladas, GWA con la velocidad va en el sentido extrapolado en 16 celdas y en contra en una, la misma ResNet-18 sobre Tiny-ImageNet con SGD. Es la lectura en velocidad de esa misma relación con la validación. GWA con el \emph{gap} se parte por arquitectura, con seis celdas a cero en la ResNet-18 en el sentido extrapolado y cero a cuatro en la convolucional, como recogió la Sección~\ref{sec:res-familia}. GSNR con la velocidad da siete celdas contra seis, con la convolucional en contra en cinco y la ResNet-18 a favor en seis, y con el \emph{accuracy} de test siete contra cuatro. La \emph{gradient disparity} con la velocidad va en contra de la extrapolación, con cuatro celdas frente a nueve. Una relación que un artículo sugiere y que no se registró como predicción se reporta solo como observación. La \emph{stiffness}, cuyo artículo la presenta como una perspectiva sobre la generalización, muestra asociación negativa con el \emph{gap} en 15 celdas y positiva en una, el sentido que ese título sugiere. No entra en el veredicto, porque una predicción registrada después de ver el signo no contrasta nada.

H6 cae en parte. De las siete predicciones, dos se cumplen en las tres arquitecturas, GWA con el \emph{accuracy} de test y GSNR con el \emph{gap}. Tres no reciben apoyo, la \emph{gradient confusion} y la escala de ruido con la velocidad y la m-coherencia con el \emph{gap}, con entre cinco y siete celdas a favor de 24 y casi tantas en contra. Dos van sistemáticamente en contra, la \emph{stiffness} con la velocidad, con nueve celdas frente a cinco, y la \emph{gradient disparity} con el \emph{accuracy} de test, con 12 frente a cuatro, las dos con la forma que la composición de curvas anticipa. El reparto tiene lógica. Una métrica leída al 5\,\% describe el estado del entrenamiento en ese momento, y ese estado anticipa dónde acaba el entrenamiento, igual que lo anticipa el \emph{accuracy} de validación. Un artículo que fijó su signo a lo largo del tiempo o del ruido de etiquetas predice otra cosa, y dentro de la celda esa lectura se confunde con el \emph{learning rate}. Lo que H6 sostiene es, por tanto, más estrecho que su enunciado. Los signos se reproducen donde el artículo habla del desenlace y la métrica describe ese desenlace, GWA y GSNR, y se invierten donde el artículo fijó el signo sobre un eje que esta matriz no recorre. Reproducirlos no convierte a ninguna métrica en una señal que supere a la validación, según decidió la Sección~\ref{sec:res-incremental}.

\section{Conclusiones del capítulo}
\label{sec:res-conclusiones}

Los datos sirven para la pregunta, con cinco salvedades medidas. La censura de la velocidad es consecuencia del diseño, y contada en pares comparables cada celda conserva entre el 62 y el 99\,\% de los suyos. El \emph{learning rate} explica entre el 45 y el 92\,\% de la dispersión de las seis métricas dentro de la celda, según la mediana sobre las 24 celdas. En las tres lecturas del predictor de referencia explica el 95\,\% o más. Toda D lleva dentro, por tanto, esa causa común. Para la velocidad solo predicen las dos primeras ventanas, y el mejor \emph{loss} de validación deja de predecir al 25\,\% en 15 celdas, por el sobreajuste. La curva de validación cuenta lo mismo que el test, en orden y en nivel. El \emph{batch} de medición de 256 ejemplos deja al estimador de la escala de ruido cerca de su valor de referencia de 255, y al 5\,\% más de un tercio de las lecturas de los entrenamientos que aprendieron queda por encima de 200. Ahí el estimador ya no resuelve las escalas de $10^{3}$ a $10^{4}$ que McCandlish et al. reportan, según deriva la Sección~\ref{sec:def-variabilidad}, y conserva el orden dentro de la celda con menos potencia.

Los seis veredictos, sobre la familia primaria al 5\,\% y los 806 entrenamientos que aprendieron, son estos. H1 se sostiene como requisito previo, porque GWA muestra asociación del mismo signo con la velocidad en 16 celdas y con el \emph{accuracy} de test en 15, más que la posición en la rejilla, y la cuenta granulada queda vacía. H2 cae, porque ninguna métrica supera a la validación de la misma ventana con intervalo ni elige mejor el \emph{learning rate}, y GWA cuesta 0,023 de \emph{accuracy} frente a 0,013. H3 cae, porque la alineación pone la primera métrica y también la última, y el orden cambia con la arquitectura. H4 se sostiene, porque cuatro métricas no maduran y dos maduran poco, mientras la validación sube hasta 0,70. H5 no se refuta ni queda demostrada, con entre cero y dos inversiones en 12 pares. H6 cae en parte, con dos predicciones cumplidas, tres sin apoyo y dos invertidas por la composición de curvas sobre el \emph{learning rate}. Ninguna clase de coste compra más capacidad predictiva que la clase gratuita.

La Tabla~\ref{tab:aportaciones} reúne, métrica a métrica, lo que las secciones anteriores reparten por hipótesis, qué lee cada una, qué anticipa al 5\,\% y en cuántas celdas, dónde falla y cómo queda frente a su artículo y frente a la validación. Una D se traduce a pares, porque $(D + 1)/2$ es la parte de los pares comparables que el predictor ordena bien. La D de 0,40 de GWA con la velocidad son siete pares de cada diez, y la de 0,71 del \emph{accuracy} de validación, ocho y medio.

\begin{table}[htbp]
\centering
\scriptsize
\setlength{\tabcolsep}{3pt}
\caption{Qué anticipa cada métrica al 5\,\%, sobre las 24 celdas y los 806 entrenamientos que aprendieron.}
\label{tab:aportaciones}
\begin{tabular}{>{\raggedright\arraybackslash}p{1.9cm}
                >{\raggedright\arraybackslash}p{2.5cm}
                >{\raggedright\arraybackslash}p{3.5cm}
                >{\raggedright\arraybackslash}p{3.0cm}
                >{\raggedright\arraybackslash}p{2.9cm}}
\toprule
\textbf{Métrica} & \textbf{Qué lee} & \textbf{Qué anticipa, celdas a favor y en contra} & \textbf{Dónde falla} & \textbf{Frente al artículo y a la validación} \\
\midrule
GWA\newline alineación; barrido, lineal & Coseno entre el gradiente y los pesos de la última capa, un margen normalizado por ejemplo & Alta, cruce más tardío (16/1) y menos \emph{accuracy} de test (15/1), en los cuatro conjuntos y las tres arquitecturas. Con el \emph{gap}, nada (6/4) & Una celda al revés, la ResNet-18 sobre Tiny-ImageNet con SGD. Con el \emph{gap} cambia de signo entre la convolucional y la ResNet-18 & Signo del artículo, 15/1. Pierde con la validación con intervalo en ocho celdas, sin ganar en ninguna; elegir con ella cuesta 0,023 frente a 0,013 \\
\addlinespace[2pt]
\emph{Stiffness}\newline alineación; barrido, cuadrática & Coseno medio entre gradientes de ejemplos de la misma clase & Alta, \emph{gap} menor (15/1), tantas celdas como la posición en la rejilla (15/0), y cruce más tardío (9/5) & La ResNet-18 se parte con el \emph{gap}; con 100 y 200 clases descansa en menos del 1\,\% de los pares; con Adam dobla su asociación con el \emph{gap} & Al revés que su artículo en velocidad (5/9). Pierde con el \emph{loss} de validación 7/17; elegir con ella cuesta 0,110, más que el azar \\
\addlinespace[2pt]
\emph{Gradient disparity}\newline variabilidad; cinco pasadas propias & Distancia entre los gradientes de dos submuestras, la raíz de la traza de la covarianza del ruido & Alta, más \emph{accuracy} de test (12/4) y \emph{gap} menor (11/4); segunda en las tres variables & Se da la vuelta en MNIST y en la ResNet-18; con Adam ordena la velocidad la mitad de fuerte que con SGD; la que más invierte entre optimizadores, cinco inversiones de 36 & Al revés que su artículo con el test (4/12). Gana a la validación con intervalo en tres celdas y pierde en ocho; cuesta 0,097 con el signo del artículo \\
\addlinespace[2pt]
GSNR\newline variabilidad; barrido, lineal & Razón señal-ruido del gradiente, parámetro a parámetro & Alto, \emph{gap} menor (11/2), en las tres arquitecturas & Con la velocidad y el test cambia de bando entre la convolucional, en contra, y la ResNet-18, a favor & Signo del artículo con el \emph{gap}, 11/2. Pierde con el \emph{loss} de validación 5/19; cuesta 0,035, la segunda mejor \\
\addlinespace[2pt]
Escala de ruido\newline variabilidad; barrido, lineal & Traza de la covarianza sobre la norma del gradiente medio, la m-coherencia con otro nombre & Nada con claridad, velocidad 7/3, test 3/6 y \emph{gap} 5/6 & La sostiene solo la ResNet-18 en velocidad, con las tres celdas en contra en MNIST; contra el techo del estimador en un tercio de las lecturas & Su artículo, 7/3 sin apoyo claro; la m-coherencia con el \emph{gap}, 5/6. Pierde con la validación 6/18; cuesta 0,073, más que el azar \\
\addlinespace[2pt]
\emph{Gradient confusion}\newline alineación; barrido, cuadrática & El coseno más negativo entre pares de ejemplos & Nada con claridad, velocidad 5/4, test 5/5 y \emph{gap} 9/5; la que más crece al esperar, de 0,18 a 0,29 con el test & Cambia de bando con la velocidad entre la red plana (4/0) y la convolucional (0/3) & Su artículo, 5/4 sin apoyo. Pierde con la validación 6/17; cuesta 0,059 \\
\bottomrule
\end{tabular}
\end{table}
